\ficha{Gambi (2015), Monopólio de emissão no segundo Banco do Brasil}{gambi2015}

\begin{identificacao}
GAMBI, Thiago Fontelas Rosado. O debate político e o pensamento econômico no
Império brasileiro: centralização de poder e monopólio de emissão no segundo
Banco do Brasil (1852-1853). \emph{Almanack}, Guarulhos, n. 9, p. 176-189,
abr. 2015. DOI 10.1590/2236-463320150913.

Artigo de história do pensamento econômico e história institucional, 14
páginas, em português. Lido integralmente na versão publicada.
\end{identificacao}

\bloco{Problema e tese}
O artigo pergunta como a controvérsia entre papelistas e metalistas guiou a
instituição do monopólio de emissão no Império, materializado no segundo Banco
do Brasil de 1853. A resposta articula economia e política. As categorias
vieram das escolas inglesas, mas foram "tropicalizadas", isto é, reprocessadas
para uma economia agrária, mercantil e escravista. E a escolha entre unidade e
pluralidade de emissão não era só técnica: seu reverso era a distribuição do
poder político entre a Corte e as províncias, de modo que a política monetária
metalista de Rodrigues Torres funcionou como braço financeiro do projeto
centralizador do Partido Conservador, sem que isso reduza sua convicção
doutrinária a mero pretexto.

\bloco{Argumento}
\begin{enumerate}[leftmargin=*,nosep]
\item O princípio bancário, da escola bancária, sustenta que a nota bancária só
se emite contra operação de crédito com contrapartida em transação comercial já
realizada, de modo que a causalidade vai do preço para a emissão e a nota
bancária não gera inflação (p. 178). O princípio monetário, ligado à teoria
quantitativa, inverte a causalidade, faz a oferta monetária determinar o preço
e nega qualquer razão para distinguir nota bancária das demais moedas (p. 179).

\item As duas escolas convergiam na conversibilidade e divergiam em dois
pontos: o direito de emissão, unidade ou pluralidade e emissor público ou
privado; e o tipo de lastro, exclusivamente ouro ou também notas do tesouro e
ações de companhias (p. 179). Na Inglaterra prevaleceu a escola monetária, com
a regra Palmer de 1827 e a lei de Peel de 1844, e a escola bancária só venceu
em crises, 1825, 1836, 1839 e 1847 (p. 180).

\item A adaptação nacional é assimétrica. Gambi retoma Saes para dizer que o
pensamento monetário europeu chegava ao Brasil, \chave{mas já era processado de
modo a refletir os problemas particulares da economia brasileira}{p. 180}
Discutir pluralidade numa economia que se industrializava era uma coisa,
discutir o mesmo tema numa economia agrária, mercantil e escravista era outra
(p. 178). O padrão-ouro foi miragem por quase todo o Império, e a paridade de 27
pence fixada em 1846 só vigorou entre 1854 e 1864, em períodos breves e sem
lastro metálico integral (p. 177).

\item Aqui os papelistas defendem a pluralidade de bancos emissores privados,
por dois motivos operacionais: garantir o giro dos negócios diante da escassez
crônica de metal, e capilarizar as notas num território vasto que um emissor
único não alcança (p. 181, p. 184). Aceitavam lastro parcial, em títulos da
dívida pública quando faltasse metal (p. 182). Sua variável de ajuste é a taxa
de juros. Os metalistas querem emissor único sob tutela do Estado, para
subordinar a provisão de liquidez à manutenção do valor da moeda; sua variável
de ajuste é a taxa de câmbio (p. 181, p. 184). Gambi sintetiza a oposição como
os homens do crédito contra os homens da moeda, ressalvando que os metalistas
do Império, na prática, cediam à realidade e mantinham o giro do comércio à
custa do lastro (p. 181, p. 184).

\item A chave política é explícita: a outra face da discussão sobre pluralidade
ou unidade de emissão é a descentralização ou a centralização do poder político,
traduzida no conflito entre Corte e províncias (p. 182-183). A correspondência
sociológica é frouxa e o autor a apresenta como tal: a unidade era defendida por
grupos do comércio externo e por proprietários escravistas do centro-sul e
também do norte, e a pluralidade por grupos ligados ao comércio interno, que
sentiam mais a escassez de moeda (p. 183).

\item O desfecho institucional. No início da década de 1850 Holanda Cavalcanti
apresentou projeto que criava bancos de emissão nas províncias, ratificando a
pluralidade então existente; Rodrigues Torres, ministro da Fazenda, respondeu
com contraprojeto de banco privado sob influência estatal e com monopólio de
emissão, e a crise monetária de 1853, na competição entre o Banco Comercial do
Rio de Janeiro e o Banco do Brasil de Irineu Evangelista de Souza, ajudou a
proposta metalista a prevalecer (p. 183). A pluralidade voltou em 1857 com
Souza Franco no Ministério da Fazenda, foi desfeita pela lei dos entraves de
1860, que sem instituir formalmente a unidade tornou a emissão inviável, e a
unidade só se restabeleceu em 1866, então como monopólio do tesouro (p. 183).

\item A seção final mostra que a divergência entre Souza Franco e Rodrigues
Torres em 1850 não era sobre a utilidade dos bancos emissores, e sim sobre a
conjuntura. Rodrigues Torres afirmava na Câmara, em 8 de março de 1850,
\chave{nas circunstâncias em que se acha o país, quando a nossa circulação é
papel-moeda, não posso ser partidista dos bancos de emissão}{p. 185} e prometia
que viria a época em que seriam úteis. Gambi lê a criação de 1853 como execução
dessa condicional, não como conversão. Souza Franco imputou ao ministro outro
motivo, \chave{todo o receio do Sr. ministro a respeito dos bancos não é porque
quebrem ou porque não vivifiquem a indústria, mas sim que o seu papel venha a
embaraçar o papel do governo: eis aqui toda a questão}{p. 187} O deputado por
São Paulo Francisco de Paula Sousa e Melo, na sessão de 26 de abril de 1850, foi
quem mais se aproximou do desenho final, um banco autorizado a ampliar emissões
na proporção do papel do governo retirado de circulação (p. 188).
\end{enumerate}

\bloco{Conceitos e definições operacionais}
Princípio bancário e princípio monetário são definidos pelo sentido da
causalidade entre transação e emissão (p. 178-179). Papelismo e metalismo são
tradução nacional desses princípios, e o autor adverte que são rótulos impostos
pela historiografia, não autodenominações dos debatedores (p. 180). Direito de
emissão: unidade ou pluralidade, mais o caráter público ou privado do emissor,
com o público justificado pelas consequências macroeconômicas da emissão e o
privado pelo lucro de emprestar notas a baixo custo (p. 179). Conversibilidade:
não se discute se, discute-se contra o quê, ouro apenas ou também notas do
tesouro e títulos (p. 179, p. 182). Nota bancária contra papel-moeda: distinção
que os metalistas recusam e que sustenta o receio de Rodrigues Torres de que a
nota do banco vire papel-moeda (p. 185-186). Bancos de depósito e desconto não
criam moeda na concepção metalista, o que explica Torres Homem autorizá-los em
1859 enquanto contraía o meio circulante (p. 185).

\bloco{Evidência e método}
História do pensamento econômico apoiada em fontes primárias parlamentares e
administrativas, sem tratamento quantitativo. As fontes centrais são os Anais
da Câmara dos Deputados nas sessões de 7 e 8 de março e 25 e 26 de abril de
1850, sobre a provincialização do meio circulante, aprovada como lei 552 de 31
de maio de 1850 e nunca executada; os relatórios do Ministério da Fazenda de
1849, 1851 e 1859; e o opúsculo \emph{Os bancos do Brasil}, de Souza Franco. A
camada interpretativa vem de Villela, Gremaud, Saes, Teixeira, Sáez, Levy e
Andrade, Guimarães, Peláez e Suzigan. O recorte é 1848 a 1853, com incursões
até 1866. O autor faz crítica explícita de fonte: é preciso saber quem fez o
discurso, com que finalidade e a que público se destinava (p. 187).

\bloco{Agentes e vertentes}
Escola bancária: Thornton, Tooke e Fullarton. Escola monetária: Ricardo, Peel e
McCulloch (p. 178-179). Antes de 1819 as mesmas posições se chamavam
bulionistas e anti-bulionistas (p. 180). No Brasil, de tendência papelista:
Bernardo de Souza Franco, deputado paraense e depois ministro da Fazenda em
1857; Irineu Evangelista de Souza, futuro diretor do banco; Francisco de Paula
Santos; José Pedro Dias de Carvalho; e Holanda Cavalcanti, autor do projeto de
bancos emissores provinciais. Villela distingue dentro do grupo o papelista
comum, como Souza Franco, que defende a conversibilidade, e o verdadeiro, como
Mauá, que prefere a moeda fiduciária (p. 182). De tendência metalista:
Rodrigues Torres, futuro presidente do banco; Ângelo Muniz da Silva Ferraz;
Francisco de Salles Torres Homem; Jerônimo José Teixeira Júnior, visconde de
Cruzeiro; e José Machado Coelho de Castro. Francisco de Paula Sousa e Melo
aparece sem filiação declarada, como formulador do arranjo que venceu. Gambi
insiste que os campos não eram homogêneos e que os homens do banco eram também
homens do governo (p. 182).

\bloco{Críticas e limites}
O artigo é de síntese e não apresenta pesquisa serial nem contrafactual; as
sessões de 1850 são lidas em passagens selecionadas. A atribuição de Thornton à
escola bancária diverge de Fonseca e Mollo (2012), que o situam entre os
bulionistas moderados, e a ficha não deve tratar a lista de escolas como
consensual. A correspondência entre corrente monetária e fração de classe é
introduzida duas vezes com o advérbio "aparentemente" (p. 180, p. 181) e não é
demonstrada. O próprio autor reconhece a arbitrariedade dos rótulos e a
dificuldade de enquadrar atores, e o exemplo que dá é forte: a experiência de
pluralidade entre o fim da década de 1830 e 1853 correu sob ministros metalistas
(p. 182). Nada no texto trata da República, da Caixa de Conversão ou do período
de 1906 a 1914, e nenhuma afirmação sobre esse período pode ser apoiada nele.

\begin{dialogo}
\item Monopólio de emissão em 1906. Procurar nos quatro jornais e nos
Documentos Parlamentares o argumento de que a Caixa concorre com o Tesouro ou
com o Banco do Brasil na função emissora, e com que vocabulário a competência
de cada um é delimitada. Gambi mostra que, em 1850, o receio decisivo do
ministro era a nota do banco virar papel-moeda e embaraçar o papel do governo.
Se esse mesmo receio reaparecer em 1906, o problema herdado é de emissor único,
não de lastro; se estiver ausente, a novidade da Caixa é justamente ter
resolvido a questão institucional pelo lastro integral.

\item Lastro contra pluralidade. Verificar se algum publicista de 1906 a 1914
propõe lastro misto para a Caixa, com títulos ou letras entrando ao lado do
ouro, que é a concessão papelista clássica descrita por Gambi (p. 182). A Caixa
opera com lastro integral em ouro. Encontrar defesa de lastro parcial mede
quanto do repertório papelista sobreviveu; não encontrar sustenta a hipótese do
capítulo de que o eixo de 1906 é o nível da taxa, não a natureza do lastro.

\item Emissão e federação. Gambi liga unidade de emissão a centralização na
Corte (p. 182-183). Procurar em 1906 se a Caixa é atacada ou defendida como
instrumento de poder de São Paulo ou da União, sobretudo no Correio Paulistano
contra os diários cariocas. Se o eixo federativo reaparecer deslocado, de Corte
contra províncias para União contra estados cafeeiros, o capítulo pode mostrar
continuidade do problema político com mudança completa do desenho monetário.

\item Homens do crédito e homens da moeda. Procurar a autodescrição pragmática,
a alegação de que a doutrina metálica é inaplicável a um país novo e agrário,
que Gambi identifica como marca da adaptação nacional. Testa se o pragmatismo
antidoutrinário permanece como recurso retórico dos defensores da Caixa, ou se
estes reivindicam, ao contrário, respaldo na ortodoxia internacional.
\end{dialogo}

\usonocapitulo{Parte III, logo depois de Fonseca e Mollo (2012). Serve a três
afirmações. Primeira, que a importação das categorias inglesas foi seletiva e
não transposição, com o critério de adaptação explicitado pelo próprio autor:
economia agrária, mercantil e escravista, escassez crônica de metal (p. 178,
p. 180). Segunda, que o tema "pluralidade de emissão contra lastro metálico" tem
data e desfecho institucional no Brasil, o embate entre o projeto de Holanda
Cavalcanti e o contraprojeto de Rodrigues Torres em 1852-1853, com a sequência
1857, 1860 e 1866 (p. 183), o que dá ao capítulo o antecedente concreto de que a
Caixa de Conversão é herdeira remota. Terceira, que a questão emissora sempre
foi também questão de distribuição de poder político, tese que o capítulo
retoma ao tratar da relação entre a Caixa, o Tesouro e os estados. O artigo não
alcança 1906, e sua função é fornecer o antecedente e o vocabulário, não
evidência sobre a Caixa.}
